\section{Experiments}
\label{experiments}

\subsection{Research Questions}
The empirical study evaluates whether the clarification workflow improves elicitation outcomes.
\begin{enumerate}[label=\textbf{RQ\arabic*:}, leftmargin=*]
  \item Does multi-agent clarification improve requirement completeness compared with non-iterative elicitation?
  \item How accurately does missing information detection identify functional, constraint, role, data, and exception gaps?
  \item How much do knowledge assistance and agent collaboration contribute to clarification effectiveness?
  \item Does the workflow reduce human effort while preserving stakeholder control?
\end{enumerate}

\subsection{Experimental Setup}
The setup compares elicitation sessions that start from the same incomplete natural-language requirements.
Each session produces clarification questions, stakeholder answers, refined requirements, and an elicitation log for evaluation.

\subsection{Datasets}
Datasets should consist of requirement elicitation scenarios with incomplete initial descriptions and reference missing-information annotations.
The dataset description should report domain coverage, scenario size, and annotation procedure without introducing results.

\subsection{Baselines}
Baselines should remain within requirement elicitation.
Candidate baselines include single-agent LLM clarification, prompt-only clarification, checklist-based elicitation, and analyst-authored clarification questions.

\subsection{Evaluation Metrics}

\subsubsection{Requirement Completeness}
Requirement completeness measures the proportion of required elicitation elements captured after clarification.

\subsubsection{Missing Information Coverage}
Missing information coverage measures how many annotated missing elements are detected and addressed by clarification questions.

\subsubsection{Clarification Effectiveness}
Clarification effectiveness measures whether generated questions lead to useful stakeholder answers and improved requirement statements.

\subsubsection{Human Effort Reduction}
Human effort reduction measures the number of clarification turns, manual edits, or reviewer interventions needed to reach an acceptable requirement.

\subsection{Main Results}
This subsection should report aggregate comparisons after experiments are completed.
No numerical claims or result tables are included before empirical execution.

\subsection{Ablation Study}
The ablation study should remove or vary major workflow components, including missing-information detection, requirement knowledge, agent collaboration, and human validation.

\subsection{Case Study}
The case study should present one requirement elicitation scenario, the detected missing information, the generated clarification questions, and the final clarified requirement statements.
